\documentclass[main]{subfiles}


\title[AutoML: Introduction]{AutoML Motivation}
\subtitle{When to use AutoML?}
\author[Marius Lindauer]{Marius Lindauer}
\date{}

\titlebackground*{images/AutoML_AutoDL}

%\institute{}
%\date{}

\begin{document}
	
	\maketitle
	

%----------------------------------------------------------------------
%----------------------------------------------------------------------
\begin{frame}[c]{You can measure it -- AutoML can optimize it!}


\begin{columns}

\column{0.5\textwidth}

\textcolor{ForestGreen}{AutoML can optimize for}:

\begin{itemize}
    \item accuracy and related metrics
    \item training and inference time
    \item memory consumption
    \item energy consumption
    \item sparsity / interpretability
    \item fairness metrics
    \item \ldots
\end{itemize}

\column{0.5\textwidth}

\pause

\textcolor{Red}{AutoML can't (directly) optimize for}

\begin{itemize}
    \item trust of users / user experience
    \item fairness in general
    \item human-driven feedback (e.g., aesthetics of images)
    \item business impact (e.g., revenue increase)
    \item emerging properties (e.g., vulnerability to unknown adversarial attacks)
    \item ...
\end{itemize}

\end{columns}

\end{frame}
%-----------------------------------------------------------------------
%----------------------------------------------------------------------
\begin{frame}[c]{Large Space of Possible Models or Pipeline Compoments}

\textcolor{ForestGreen}{Do Use AutoML}
\begin{itemize}
    \item Humans are fairly bad at navigating huge design spaces if there is no intuition
    \item[$\leadsto$] Use AutoML to quickly prototype a first solution
\end{itemize}

\pause

\textcolor{Red}{Don't Use AutoML}

\begin{itemize}
    \item If there are already known constraints on the ML model, you should consider these first and not blindly apply AutoML
    \item[$\leadsto$] Constraints might directly imply a particular (type of) model
\end{itemize}

\end{frame}
%-----------------------------------------------------------------------
%----------------------------------------------------------------------
\begin{frame}[c]{No Good Baseline and Missing Expertise}

\textcolor{ForestGreen}{Do Use AutoML}
\begin{itemize}
    \item Sometimes, even experienced ML developers lack an intuition of what to try out first if they are confronted with a completely new task
    \item Non-technical users (e.g., domain experts) are facing this even more often
    \item[\leadsto] AutoML can help you with identifying promising directions and exclude non-working approaches $\to$ AutoML for Rapid Prototyping
\end{itemize}

\pause

\textcolor{Red}{Don't Use AutoML}

\begin{itemize}
    \item If you already know a promising baseline (e.g., on similar datasets), you should try these out first
    \item AutoML tools often need more compute compared to a single model training
    \begin{itemize}
        \item However, if you tried $10$ different models, it could be worth considering AutoML instead!
    \end{itemize}
\end{itemize}

\end{frame}
%-----------------------------------------------------------------------
%----------------------------------------------------------------------
\begin{frame}[c]{Well-performing Hyperparameters Settings are Unknown}

\textcolor{ForestGreen}{Do Use AutoML}
\begin{itemize}
    \item Beyond trying out default or well-established settings of hyperparameters, humans are fairly bad at efficiently trying out different hyperparameters
    \begin{itemize}
        \item That is even more so in high-dimensional spaces and if hyperparameters interact with each other
    \end{itemize}
    %\item Exhaustive (grid) search is often very expensive in terms of computing
    %\item[$\leadsto$] Bayesian Optimization is a well-established and more efficient method for tuning hyperparameters 
\end{itemize}

\pause

\textcolor{Red}{Don't Use AutoML}

\begin{itemize}
    \item If you deal with a well-known dataset (e.g., Imagenet), there are often very well-performing hyperparameter configurations for all kinds of things
    \item[$\leadsto$] Use these known hyperparameter settings  at least as the very first configuration your AutoML tries out
\end{itemize}

\end{frame}
%-----------------------------------------------------------------------
%----------------------------------------------------------------------
\begin{frame}[c]{Explainability and Documentation!?}

\textcolor{ForestGreen}{Do Use AutoML}
\begin{itemize}
    \item Results of AutoML tools (e.g., performances of different models) can be automatically saved and reproduced
    \item Such empirical results can be used by an explainability method to better understand the design process (e.g., which hyperparameter was the most important one)
\end{itemize}

\pause

\textcolor{Red}{Don't Use AutoML}

\begin{itemize}
    \item Some AutoML tools generate fairly complex ensemble/stacking models, which might be harder to interpret than simpler models with less accuracy.
\end{itemize}

\end{frame}
%-----------------------------------------------------------------------
%----------------------------------------------------------------------
\begin{frame}[c]{Fairness and AutoML \liturl{Weerts et al. 2023}{https://arxiv.org/abs/2303.08485} }

\textcolor{ForestGreen}{Do Use AutoML}
\begin{itemize}
    \item If you already know that there exists a fairness metric which is relevant for the application, you can use AutoML for setting constraints on fairness (mitigating discrimination)
\end{itemize}

\pause

\textcolor{Red}{Don't Use AutoML}

\begin{itemize}
    \item If you do not know which fairness metric to use and discrimination could be a problem for your application, you should not apply AutoML blindly 
    \begin{itemize}
        \item Not considering fairness could lead to unwanted results
        \item Considering a random fairness metric by chance could also lead to arbitrarily bad results in terms of the actual discrimination risk
    \end{itemize}
\end{itemize}

\end{frame}
%-----------------------------------------------------------------------

%----------------------------------------------------------------------
\begin{frame}{Summary} 

Most important insights from this video:

\begin{enumerate}
    \item AutoML can be very useful to increase the efficiency of developing new ML applications
    \item AutoML can do a lot of things for you, but there are certain aspects that cannot be automated (e.g.,  what fairness metric to optimize for)
    \item If you already have insights into your problem, make use of them first! 
    \begin{itemize}
        \item If you are not satisfied, you can still give it a try with AutoML
    \end{itemize}
\end{enumerate}


\end{frame}

\end{document}

